\documentclass[aspectratio=169,11pt]{beamer}
\usepackage[utf8]{inputenc}
\usepackage[T1]{fontenc}
\usepackage[spanish,es-noshorthands]{babel}
\usepackage{lmodern}
\usepackage{tikz}
\usetikzlibrary{arrows.meta,positioning,calc,shapes.geometric}
\usepackage{booktabs}
\usepackage{csquotes}
\usepackage{hyperref}
\usepackage{xurl} % permite partir URLs largas en la bibliografía
\hypersetup{colorlinks=true,linkcolor=black,urlcolor=black,citecolor=black}
\usepackage[backend=biber,style=ieee,sorting=none,maxnames=3]{biblatex}
\addbibresource{referencias.bib}

\definecolor{azul}{HTML}{123B2A}
\definecolor{azulclaro}{HTML}{EAF3EE}
\definecolor{turquesa}{HTML}{2E6B57}
\definecolor{turquesaclaro}{HTML}{D9EDE2}
\definecolor{tinta}{HTML}{1F2D27}
\definecolor{gris}{HTML}{55665D}
\definecolor{fondo}{HTML}{FBFCFB}

\setbeamercolor{normal text}{fg=tinta,bg=fondo}
% La bibliografía no hereda el verde del color estructural del tema.
\setbeamercolor{bibliography item}{fg=black}
\setbeamercolor{bibliography entry author}{fg=black}
\setbeamercolor{bibliography entry title}{fg=black}
\setbeamercolor{bibliography entry location}{fg=black}
\setbeamercolor{bibliography entry note}{fg=black}
\setbeamercolor{bibliography entry journal}{fg=black}
\AtBeginBibliography{\color{black}\sloppy\emergencystretch=2em}

\setbeamercolor{structure}{fg=turquesa}
\setbeamercolor{frametitle}{fg=azul,bg=fondo}
\setbeamercolor{block title}{fg=white,bg=azul}
\setbeamercolor{block body}{fg=tinta,bg=azulclaro}
\setbeamertemplate{navigation symbols}{}
\setbeamertemplate{itemize item}{\color{turquesa}\tiny$\blacksquare$}
\setbeamertemplate{footline}{%
  \hfill{\color{gris}\scriptsize Aplicaciones Distribuidas · UTEQ\hspace{0.7em}|\hspace{0.7em}\insertframenumber}\hspace{0.38cm}\vspace{0.13cm}}
\setbeamertemplate{frametitle}{\vspace{0.16cm}\hspace{0.12cm}\insertframetitle\par\vspace{-0.04cm}%
  \hspace{0.12cm}{\color{turquesa}\rule{0.86\paperwidth}{0.7pt}}\vspace{0.15cm}}

\tikzset{
   caja/.style={rounded corners=5pt,draw=azul!25,fill=azulclaro,minimum height=0.86cm,align=center,inner sep=5pt},
   evento/.style={rounded corners=5pt,draw=turquesa!70,fill=turquesaclaro,minimum height=0.86cm,align=center,inner sep=5pt},
   flecha/.style={-{Latex[length=2.5mm]},thick,draw=gris},
   propuesta/.style={-{Latex[length=2.2mm]},thick,dashed,draw=turquesa}
}
\newcommand{\fuentes}[1]{\vfill\hfill{\tiny\color{gris}Fuentes: #1}\vspace{-0.11cm}}

\title{Mensajería Publish/Subscribe}
\subtitle{MQTT, JMS y AMQP · Aplicación propuesta al SGRA}
\author{Grupo de exposición}
\date{Noviembre de 2026}

\begin{document}

\begin{frame}[plain]
\begin{tikzpicture}[remember picture,overlay]
  \fill[azul] (current page.south west) rectangle (current page.north east);
  \fill[turquesa] ([xshift=-0.55cm]current page.north east) -- ([xshift=-2.7cm]current page.north east)
    -- ([xshift=-0.55cm,yshift=-2.2cm]current page.north east) -- cycle;
  \draw[line width=1.4pt,turquesa!75] ([xshift=0.72cm,yshift=0.22cm]current page.south west) -- ([xshift=12.2cm,yshift=0.22cm]current page.south west);
\end{tikzpicture}
\vspace{0.20cm}
{\color{white}\small UNIVERSIDAD TÉCNICA ESTATAL DE QUEVEDO}\\[1mm]
{\color{white!75}\scriptsize Facultad de Ciencias de la Computación · Carrera de Software}\\[12mm]
{\color{white}\fontsize{24}{28}\selectfont\bfseries Mensajería Publish/Subscribe}\\[2mm]
{\color{white!88}\Large MQTT · JMS · AMQP}\\[2mm]
{\color{white!75}\normalsize Aplicaciones propuestas a SGRA, CLAB, RTV y SSDC}
\vspace{0.6cm}
\begin{columns}[onlytextwidth,T]
\column{0.49\textwidth}
{\color{white}\small\textbf{Integrantes}}\\[2mm]
{\color{white!87}\scriptsize Zambrano Parraga Luis Mario}\\[1.2mm]
{\color{white!87}\scriptsize Manotoa Alava Carlos Alexander}\\[1.2mm]
{\color{white!87}\scriptsize Piguave Montaño Nixon Antonio}\\[1.2mm]
{\color{white!87}\scriptsize Silva Triviño John Jairo}
\column{0.48\textwidth}
{\color{white}\small\textbf{Asignatura}}\\[1.6mm]
{\color{white!87}\scriptsize Aplicaciones Distribuidas · ISR-701}\\[1mm]
{\color{white!87}\scriptsize Séptimo nivel · Paralelo A}\\[1mm]
{\color{white!87}\scriptsize Docente: Gleiston C. Guerrero-Ulloa, PhD}
\end{columns}
\end{frame}

% Cada integrante edita exclusivamente su archivo de seccion.
% Sección inicial: Mario Zambrano.
% Las afirmaciones académicas llevan \cite{} en la propia diapositiva;
% la arquitectura existente de SGRA se documenta mediante sus repositorios publicos.

\begin{frame}{SGRA: punto de partida}
\begin{columns}[T,onlytextwidth]
\column{0.48\textwidth}
{\small\bfseries Arquitectura documentada}\vspace{2mm}
\begin{center}
\begin{tikzpicture}[node distance=0.45cm,scale=0.88,transform shape]
\node[caja,minimum width=4.15cm] (ui) {\textbf{Angular 21}\\\scriptsize Interfaz web};
\node[caja,minimum width=4.15cm,below=of ui] (api) {\textbf{Spring Boot}\\\scriptsize API REST + reglas};
\node[caja,minimum width=4.15cm,below=of api] (db) {\textbf{PostgreSQL}\\\scriptsize Datos de tutorías};
\draw[flecha] (ui) -- node[right,font=\scriptsize]{HTTPS} (api);
\draw[flecha] (api) -- (db);
\end{tikzpicture}
\end{center}
\column{0.49\textwidth}
{\small\bfseries Necesidad planteada}\vspace{3mm}
\begin{tikzpicture}
\node[evento,text width=4.8cm,minimum height=1.1cm] {\bfseries Se registra una solicitud\\de tutoría};
\end{tikzpicture}
\vspace{3mm}
\begin{itemize}\setlength{\itemsep}{4pt}
  \item Notificar al docente
  \item Registrar auditoría
  \item Actualizar indicadores
\end{itemize}
\vspace{2mm}
{\scriptsize\color{turquesa}\textbf{Pregunta:} ¿cómo comunicar el mismo evento a varios componentes?}
\end{columns}
\fuentes{\cite{sgra_frontend,sgra_backend,microsoft2026}}
\end{frame}

\begin{frame}{PUB/SUB: un evento, varios receptores}
\begin{center}
\begin{tikzpicture}[scale=0.94,transform shape]
\node[caja,minimum width=2.55cm,minimum height=1.3cm] (publisher) at (0,0) {\textbf{Publicador}\\\scriptsize SGRA / backend};
\node[evento,minimum width=2.8cm,minimum height=1.3cm] (broker) at (4.3,0) {\textbf{Broker}\\\scriptsize tópico: tutorías};
\node[caja,minimum width=2.55cm] (s1) at (8.8,1.35) {\textbf{Suscriptor 1}\\\scriptsize Notificaciones};
\node[caja,minimum width=2.55cm] (s2) at (8.8,0) {\textbf{Suscriptor 2}\\\scriptsize Auditoría};
\node[caja,minimum width=2.55cm] (s3) at (8.8,-1.35) {\textbf{Suscriptor 3}\\\scriptsize Indicadores};
\draw[flecha] (publisher) -- node[above,font=\scriptsize]{publica evento} (broker);
\foreach \n in {s1,s2,s3}{\draw[flecha] (broker.east) -- (\n.west);}
\end{tikzpicture}
\end{center}
\vspace{0.15cm}
\begin{columns}[T,onlytextwidth]
\column{0.32\textwidth}\centering\small\bfseries Desacoplamiento
\column{0.35\textwidth}\centering\small\bfseries Comunicación asíncrona
\column{0.31\textwidth}\centering\small\bfseries Entrega por suscripción
\end{columns}
\fuentes{\cite{lazidis2022,dong2024,kamoun2024,microsoft2026}}
\end{frame}

\begin{frame}{Mini práctica local: un evento, dos suscriptores}
{\small\color{gris}Simulación de \texttt{SolicitudTutoriaCreada}; no es una función ya implementada en SGRA.}
\vspace{2.5mm}
\begin{columns}[T,onlytextwidth]
\column{0.49\textwidth}
\begin{block}{Terminal A · Notificaciones}
{\footnotesize\ttfamily mosquitto\_sub -t\\ sgra/tutorias/creada -v}
\end{block}
\column{0.49\textwidth}
\begin{block}{Terminal B · Auditoría}
{\footnotesize\ttfamily mosquitto\_sub -t\\ sgra/tutorias/creada -v}
\end{block}
\end{columns}
\vspace{3mm}
\begin{block}{Terminal C · Publicación del evento}
{\footnotesize\ttfamily mosquitto\_pub -t sgra/tutorias/creada\\ -m "Solicitud 42 creada"}
\end{block}
\begin{center}
{\small\color{turquesa}\bfseries Resultado esperado: ambos suscriptores reciben el mismo mensaje.}
\end{center}
\fuentes{\cite{mosquitto_pub,mosquitto_sub,mqtt2019}}
\end{frame}

\begin{frame}{SGRA: evolución orientada a eventos}
\begin{center}
\begin{tikzpicture}[scale=0.84,transform shape]
\node[caja,minimum width=2.45cm] (angular) at (0,2.2) {\textbf{Angular}\\\scriptsize Interfaz};
\node[caja,minimum width=3cm] (api) at (4.2,2.2) {\textbf{Spring Boot}\\\scriptsize API y negocio};
\node[caja,minimum width=2.4cm] (postgres) at (8.7,2.2) {\textbf{PostgreSQL}\\\scriptsize Persistencia};
\draw[flecha] (angular) -- node[above,font=\scriptsize]{REST/HTTPS} (api);
\draw[flecha] (api) -- (postgres);
\node[evento,minimum width=3cm] (broker) at (4.2,-0.1) {\textbf{Broker de eventos}\\\scriptsize Componente propuesto};
\draw[propuesta] (api) -- node[right,font=\scriptsize]{publicación} (broker);
\node[caja,minimum width=2.65cm] (mail) at (0.1,-2.0) {\textbf{Notificaciones}};
\node[caja,minimum width=2.65cm] (audit) at (4.2,-2.0) {\textbf{Auditoría}};
\node[caja,minimum width=2.65cm] (stats) at (8.7,-2.0) {\textbf{Indicadores}};
\foreach \n in {mail,audit,stats}{\draw[propuesta] (broker.south) -- (\n.north);}
\end{tikzpicture}
\end{center}
\vspace{-0.05cm}
{\scriptsize\textbf{Ventaja:} nuevas funciones sin llamadas directas desde el productor.\quad
\textbf{Reto:} entrega fiable, errores y consistencia entre componentes.}
\fuentes{\cite{sgra_frontend,sgra_backend,microsoft2026,maharjan2023,derlatka2023,bezerra2024}}
\end{frame}

% Carlos Manotoa: MQTT. Paleta, cabecera y pie heredados de main.tex.
% Autocontenido: gráficos vectoriales TikZ, sin imágenes externas.
% El texto de exposición se conserva en notas Beamer, ocultas por defecto.
\begingroup
\section{Protocolo MQTT}

% Ilustraciones vectoriales editables, sin archivos externos.
\newcommand{\mqttIcon}[3]{%
\begingroup
\pgfmathtruncatemacro{\mqttIconId}{#1+3*#2}%
\resizebox{#3}{#3}{%
\begin{tikzpicture}[x=1cm,y=1cm,line width=.9pt,
  line cap=round,line join=round,draw=azul]
\path[use as bounding box] (0,0) rectangle (4,4);
\clip (0,0) rectangle (4,4);
\ifcase\mqttIconId
% 0: Portátil.
\shade[top color=white,bottom color=turquesaclaro,draw=azul,rounded corners=3pt]
 (.55,1.55)--(.38,3.43)--(2.94,3.43)--(3.15,1.55)--cycle;
\path[fill=white,draw=turquesa] (.73,1.79)--(.59,3.23)--(2.76,3.23)--(2.92,1.79)--cycle;
\draw[turquesa,line width=1.4pt] (1.14,2.16)--(1.72,2.83)--(2.35,2.16);
\foreach \x/\y in {1.14/2.16,1.72/2.83,2.35/2.16}{\draw[fill=turquesaclaro] (\x,\y) circle (.16);}
\shade[top color=azulclaro,bottom color=turquesaclaro,draw=azul]
 (.55,1.55)--(3.15,1.55)--(3.75,.82)--(.83,.82)--cycle;
\path[fill=turquesa!35,draw=azul] (.83,.82)--(3.75,.82)--(3.64,.68)--(.85,.68)--cycle;
\foreach \y/\xl/\xr in {1.37/.83/3.04,1.19/.94/3.20}{\draw[gris,line width=.55pt] (\xl,\y)--(\xr,\y);}
\foreach \x in {1.05,1.36,1.67,1.98,2.29,2.60,2.91}{\draw[gris,line width=.45pt] (\x,1.48)--(\x+.13,1.08);}
\draw[turquesa,fill=white] (1.65,1.02)--(2.52,1.02)--(2.66,.88)--(1.77,.88)--cycle;
\or
% 1: Servidor.
\path[fill=white,draw=turquesa!65] (.72,2.73)
 .. controls (.18,2.75) and (.12,3.34) .. (.8,3.39)
 .. controls (.78,3.84) and (1.48,4.03) .. (1.83,3.6)
 .. controls (2.2,4.1) and (2.98,3.9) .. (3.05,3.41)
 .. controls (3.87,3.5) and (3.83,2.74) .. (3.16,2.73)--cycle;
\path[fill=turquesa!55,draw=azul] (2.77,.6)--(3.29,.91)--(3.29,3.22)--(2.77,2.95)--cycle;
\shade[left color=white,right color=turquesaclaro,draw=azul]
 (.85,2.95)--(1.37,3.22)--(3.29,3.22)--(2.77,2.95)--cycle;
\shade[top color=turquesaclaro,bottom color=turquesa!40,draw=azul,rounded corners=2pt]
 (.85,.6) rectangle (2.77,2.95);
\foreach \y in {.85,1.52,2.19}{
 \path[fill=azulclaro,draw=azul,rounded corners=2pt] (1.02,\y) rectangle (2.59,\y+.47);
 \draw[fill=turquesa] (1.24,\y+.235) circle (.065);
 \draw[turquesa,line width=1pt] (1.45,\y+.235)--(1.85,\y+.235);
 \foreach \x in {2.10,2.24,2.38}{\draw[azul,line width=1pt] (\x,\y+.12)--(\x,\y+.35);}
}
\path[fill=azul] (1,.45) rectangle (1.42,.60);
\path[fill=azul] (2.18,.45) rectangle (2.6,.60);
\or
% 2: Notificaciones.
\draw[fill=turquesaclaro] (2,3.53) circle (.12);
\shade[left color=white,right color=turquesa!40,draw=azul]
 (1.13,2.15) .. controls (1.55,2.4) and (1.15,3.44) .. (2,3.44)
 .. controls (2.85,3.44) and (2.45,2.4) .. (2.87,2.15)--cycle;
\draw[fill=turquesaclaro] (1.98,2.13) ellipse (.94 and .12);
\path[fill=azul] (1.8,1.99) arc[start angle=180,end angle=360,radius=.2];
\draw[turquesa] (.85,2.75) arc[start angle=205,end angle=155,radius=.65];
\draw[turquesa] (3.15,2.75) arc[start angle=-25,end angle=25,radius=.65];
\path[fill=white,draw=azul,rounded corners=2pt] (.47,.62) rectangle (2.78,1.75);
\draw[turquesa] (.49,1.71)--(1.63,.93)--(2.76,1.71);
\draw[turquesa] (.5,.65)--(1.22,1.14);
\draw[turquesa] (2.75,.65)--(2.03,1.14);
\draw[fill=turquesaclaro] (3.13,1.55) circle (.35);
\shade[top color=turquesaclaro,bottom color=turquesa!45,draw=azul]
 (2.55,.48)--(2.55,.76) .. controls (2.55,1.29) and (3.7,1.29)
 .. (3.7,.76)--(3.7,.48)--cycle;
\or
% 3: Auditoría.
\shade[left color=white,right color=turquesa!40,draw=azul]
 (.25,.68)--(.25,2.13) arc[start angle=180,end angle=360,x radius=.75,y radius=.22]
 --(1.75,.68) arc[start angle=0,end angle=-180,x radius=.75,y radius=.22];
\draw[fill=azulclaro] (1,2.13) ellipse (.75 and .22);
\foreach \y in {1.12,1.59}{\draw[turquesa] (.25,\y) arc[start angle=180,end angle=360,x radius=.75,y radius=.22];}
\shade[top color=turquesaclaro,bottom color=turquesa!35,draw=azul,rounded corners=4pt]
 (1.2,.42) rectangle (3.65,3.34);
\path[fill=white,draw=azul,rounded corners=2pt] (1.37,.61) rectangle (3.48,3.12);
\path[fill=turquesaclaro,draw=azul,rounded corners=2pt] (1.91,2.99) rectangle (2.94,3.46);
\draw[fill=turquesaclaro] (2.425,3.49) circle (.16);
\foreach \y in {1.02,1.75,2.48}{
 \draw[turquesa] (1.63,\y) rectangle (2.08,\y+.38);
 \draw[azul,line width=1.5pt] (1.68,\y+.18)--(1.82,\y+.05)--(2.14,\y+.43);
 \draw[turquesa] (2.38,\y+.28)--(3.15,\y+.28);
 \draw[turquesa!60] (2.38,\y+.09)--(2.96,\y+.09);
}
\or
% 4: Base de datos.
\foreach \y in {.66,1.52,2.38}{
 \shade[left color=white,right color=turquesa!45,draw=azul]
 (.55,\y)--(.55,\y+.61) arc[start angle=180,end angle=360,x radius=1.45,y radius=.38]
 --(3.45,\y) arc[start angle=0,end angle=-180,x radius=1.45,y radius=.38];
 \draw[fill=azulclaro] (2,\y+.61) ellipse (1.45 and .38);
}
\or
% 5: Estudiante, figura esquemática.
\shade[top color=turquesaclaro,bottom color=turquesa!55,draw=azul]
 (.72,.4)--(.87,1.35) .. controls (1.1,1.84) and (2.9,1.84)
 .. (3.13,1.35)--(3.28,.4)--cycle;
\path[fill=azulclaro,draw=azul] (1.72,1.55) rectangle (2.28,2.03);
\draw[fill=white] (2,2.48) ellipse (.67 and .82);
\path[fill=azul] (1.36,2.66) .. controls (1.1,3.49) and (2.65,3.61)
 .. (2.67,2.65)--(2.51,2.97)--(2.08,3.12)--(1.64,2.87)--cycle;
\fill[azul] (1.76,2.54) circle (.045);
\fill[azul] (2.24,2.54) circle (.045);
\draw[turquesa] (1.82,2.19) .. controls (2,2.07) and (2.12,2.11) .. (2.22,2.23);
\path[fill=turquesa,draw=azul] (1.36,3.18)--(1.36,3.5)--(2.64,3.5)--(2.64,3.18)
 .. controls (2.25,3.04) and (1.75,3.04) .. (1.36,3.18);
\path[fill=turquesaclaro,draw=azul] (.9,3.5)--(2,3.91)--(3.12,3.5)--(2,3.1)--cycle;
\draw[azul] (2.95,3.48)--(2.95,2.95);
\draw[azul,line width=2pt] (2.95,2.96)--(2.95,2.79);
\path[fill=white,draw=azul,rounded corners=2pt] (1.33,.45) rectangle (2.63,1.48);
\path[fill=turquesaclaro,draw=azul] (1.4,.51) rectangle (2.63,1.4);
\draw[turquesa] (1.51,.6)--(1.51,1.31);
\draw[turquesa] (1.7,1.12)--(2.36,1.12);
\fi
\end{tikzpicture}}%
\endgroup
}

\tikzset{
  mqttlabel/.style={font=\small\bfseries,text=azul,align=center},
  mqttdata/.style={font=\scriptsize\ttfamily,text=tinta,align=center},
  mqttflow/.style={-{Latex[length=2.5mm]},line width=1.2pt,draw=turquesa},
  mqttsub/.style={-{Latex[length=2mm]},dashed,line width=.8pt,draw=gris},
  mqttpacket/.style={-{Latex[length=1.7mm]},line width=.7pt,draw=turquesa},
  mqttillustration/.style={inner sep=0pt,outer sep=0pt}
}
% Reserva una línea propia encima del pie; no modifica main.tex.
\renewcommand{\fuentes}[1]{%
  \begin{tikzpicture}[remember picture,overlay]
    \node[anchor=south west,font=\tiny,text=gris]
      at ([xshift=.55cm,yshift=.48cm]current page.south west)
      {Fuentes: #1};
  \end{tikzpicture}%
}

% ------------------------------------------------------------
% 1. Arquitectura y secuencias QoS
% ------------------------------------------------------------
\begin{frame}{MQTT · Arquitectura y QoS}
\begin{tikzpicture}[remember picture,overlay,shift={(current page.south west)}]
  \node[mqttillustration] (pub) at (2.0,5.7) {\mqttIcon{0}{0}{2.65cm}};
  \node[mqttillustration] (broker) at (7.5,5.7) {\mqttIcon{1}{0}{2.75cm}};
  \node[mqttillustration] (suba) at (13.2,6.65) {\mqttIcon{2}{0}{1.8cm}};
  \node[mqttillustration] (subb) at (13.2,4.65) {\mqttIcon{0}{1}{1.8cm}};
  \node[mqttlabel] at (2.0,4.15) {Publicador};
  \node[mqttlabel] at (7.5,4.15) {Broker MQTT};
  \node[mqttlabel] at (13.2,5.6) {Suscriptor A};
  \node[mqttlabel] at (13.2,3.6) {Suscriptor B};
  \node[mqttdata] at (2.0,7.35) {TCP/IP};
  \draw[mqttflow] (3.4,5.7) -- node[above,font=\scriptsize]{PUBLISH} (6.0,5.7);
  \draw[mqttflow] (9.0,5.7) -- (10.0,5.7) |- (12.2,6.65);
  \draw[mqttflow] (10.0,5.7) |- (12.2,4.65);
  \node[mqttdata] at (9.55,7.25) {sgra/tutorias/creada};
  \draw[draw=turquesa!35] (.8,3.35) -- (15.2,3.35);

  % Cada diagrama muestra un único tramo emisor--receptor MQTT.
  \foreach \x/\q in {0.75/0,5.65/1,10.55/2}{
    \node[mqttlabel,anchor=west] at (\x,3.0) {QoS \q};
    \node[font=\scriptsize,text=gris] at (\x+.65,2.58) {Emisor};
    \node[font=\scriptsize,text=gris] at (\x+3.5,2.58) {Receptor};
    \draw[draw=gris!35] (\x+.65,2.37) -- (\x+.65,.98);
    \draw[draw=gris!35] (\x+3.5,2.37) -- (\x+3.5,.98);
  }
  \draw[mqttpacket] (1.4,2.12) -- node[above,font=\scriptsize]{PUBLISH} (4.25,2.12);
  \draw[mqttpacket] (6.3,2.12) -- node[above,font=\scriptsize]{PUBLISH} (9.15,2.12);
  \draw[mqttpacket] (9.15,1.51) -- node[above,font=\scriptsize]{PUBACK} (6.3,1.51);
  \draw[mqttpacket] (11.2,2.22) -- node[above,font=\scriptsize]{PUBLISH} (14.05,2.22);
  \draw[mqttpacket] (14.05,1.84) -- node[above,font=\scriptsize]{PUBREC} (11.2,1.84);
  \draw[mqttpacket] (11.2,1.46) -- node[above,font=\scriptsize]{PUBREL} (14.05,1.46);
  \draw[mqttpacket] (14.05,1.08) -- node[above,font=\scriptsize]{PUBCOMP} (11.2,1.08);
\end{tikzpicture}
\fuentes{\cite{shahri2022mqtt,im2023emqtt,innamorati2026mqtt}}
\note{
MQTT es un protocolo ligero de mensajería de la capa de aplicación,
habitualmente sobre TCP/IP. Los clientes pueden publicar, suscribirse
o desempeñar ambos papeles. El broker recibe las publicaciones y las
distribuye según las suscripciones. El publicador no necesita conocer
la dirección ni la cantidad de receptores.
\par
El tópico de esta exposición es sgra/tutorias/creada. Los niveles se
separan mediante barras. Los filtros admiten + para un nivel y \# para
los niveles restantes; una suscripción a sgra/tutorias/\# puede recibir
varios tipos de eventos bajo esa jerarquía. El esquema superior muestra
la distribución de publicaciones; la suscripción se establece antes.
\par
Los diagramas inferiores representan un único intercambio emisor--receptor.
QoS 0 envía PUBLISH sin confirmación MQTT: entrega como máximo una vez,
con posibilidad de pérdida. QoS 1 utiliza PUBLISH y PUBACK: entrega al
menos una vez, con posibles duplicados. QoS 2 utiliza PUBLISH, PUBREC,
PUBREL y PUBCOMP para entregar exactamente una vez en ese intercambio.
\par
Publicador--broker y broker--suscriptor son intercambios distintos.
El acuse del broker no demuestra que el consumidor haya terminado su
tarea. QoS 2 no garantiza una única ejecución de la operación de negocio.
Tampoco los niveles QoS establecen un plazo máximo de entrega.
}
\end{frame}

% ------------------------------------------------------------
% 2. Demostración local. Datos esperados, no evidencia fabricada.
% ------------------------------------------------------------
\begin{frame}{Mosquitto · MQTTX}
\begin{tikzpicture}[remember picture,overlay,shift={(current page.south west)}]
  \node[mqttillustration] at (2.05,5.45) {\mqttIcon{0}{0}{2.65cm}};
  \node[mqttillustration] at (7.5,5.45) {\mqttIcon{1}{0}{2.65cm}};
  \node[mqttillustration] at (13.2,6.2) {\mqttIcon{0}{0}{1.8cm}};
  \node[mqttillustration] at (13.2,4.1) {\mqttIcon{0}{0}{1.8cm}};
  \node[mqttlabel] at (2.05,7.25) {MQTTX};
  \node[mqttlabel] at (7.5,7.25) {Mosquitto};
  \node[mqttdata] at (7.5,6.9) {127.0.0.1:1883};
  \node[mqttdata] at (2.05,3.95) {pub-sgra};
  \node[mqttlabel] at (13.2,7.25) {MQTTX};
  \node[mqttdata] at (13.2,5.15) {sub-notificaciones};
  \node[mqttdata] at (13.2,3.05) {sub-auditoria};
  \draw[mqttflow] (3.4,5.45) -- node[above,font=\scriptsize]{PUBLISH} (6.0,5.45);
  \draw[mqttflow] (9.0,5.3) -- (10.0,5.3) |- (12.2,6.0);
  \draw[mqttflow] (10.0,5.3) |- (12.2,3.9);
  \draw[mqttsub] (12.2,6.5) -- node[above,font=\scriptsize]{SUBSCRIBE} (9.0,6.5);
  \draw[mqttsub] (12.2,4.4) -- (10.5,4.4) -- (9.0,4.85);
  \node[mqttdata] at (7.5,3.85) {QoS 1\quad Retain 0};
  \node[mqttdata,font=\small\ttfamily] at (7.5,2.75) {sgra/tutorias/creada};
  \draw[draw=turquesa!35] (.8,2.3) -- (15.2,2.3);
  \node[mqttdata,text=gris] at (2.4,1.9) {eventoId};
  \node[mqttdata,text=gris] at (7.8,1.9) {tipo};
  \node[mqttdata,text=gris] at (13.3,1.9) {solicitudId};
  \node[mqttlabel] at (2.4,1.4) {evt-001};
  \node[mqttdata,font=\small\ttfamily] at (7.8,1.4) {SolicitudTutoriaCreada};
  \node[mqttlabel,font=\Large\bfseries] at (13.3,1.4) {42};
\end{tikzpicture}
\fuentes{\cite{mqtt2019,mosquitto_pub,mosquitto_sub}; \href{https://mqttx.app/docs}{MQTTX}}
\note{
Esta es la configuración de una demostración local con datos ficticios,
sin conexión con el SGRA real. Mosquitto es el broker y MQTTX Desktop
es el cliente gráfico. Las tres conexiones tienen identificadores distintos.
\par
Primero se conectan sub-notificaciones y sub-auditoria y se suscriben a
sgra/tutorias/creada solicitando QoS 1. Las flechas discontinuas indican
esas suscripciones. Después pub-sgra publica un JSON en el mismo tópico,
con QoS 1 y Retain desactivado. Las flechas continuas representan el
recorrido esperado de esa publicación.
\par
El JSON tiene tres campos: eventoId con valor evt-001; tipo con valor
SolicitudTutoriaCreada; y solicitudId con valor numérico 42.
El identificador del evento y el de la solicitud cumplen funciones diferentes.
\par
El resultado esperado es recibir el mismo evento en ambos suscriptores.
No se afirma que esta ejecución ya se haya realizado: se necesitan
capturas auténticas después de probarla. La ilustración es un esquema,
no una captura de MQTTX. Una recepción satisfactoria no demuestra
recuperación ante fallos. La configuración del ZIP no persiste mensajes en disco.
}
\end{frame}

% ------------------------------------------------------------
% 3. SGRA: arquitectura de referencia y mensajería propuesta.
% ------------------------------------------------------------
\begin{frame}{SGRA · Propuesta MQTT}
\begin{tikzpicture}[remember picture,overlay,shift={(current page.south west)}]
  \node[mqttillustration] at (1.5,6.15) {\mqttIcon{2}{1}{2.1cm}};
  \node[mqttillustration] at (5.05,6.15) {\mqttIcon{0}{0}{2.1cm}};
  \node[mqttillustration] at (8.8,6.15) {\mqttIcon{1}{0}{2.1cm}};
  \node[mqttillustration] at (13.25,6.15) {\mqttIcon{1}{1}{2.1cm}};
  \node[mqttlabel] at (1.5,4.9) {Estudiante};
  \node[mqttlabel] at (5.05,4.9) {Angular};
  \node[mqttlabel] at (8.8,4.9) {Spring Boot};
  \node[mqttlabel] at (13.25,4.9) {PostgreSQL};
  \draw[mqttflow] (2.65,6.15) -- (3.9,6.15);
  \draw[mqttflow] (6.2,6.15) -- node[above,font=\scriptsize]{REST/HTTPS} (7.65,6.15);
  \draw[mqttflow] (9.95,6.15) -- node[above,font=\scriptsize]{SQL} (12.1,6.15);

  \node[mqttillustration] at (8.8,2.95) {\mqttIcon{1}{0}{1.95cm}};
  \node[mqttlabel] at (8.8,1.8) {Mosquitto};
  \node[mqttillustration] at (13.25,3.5) {\mqttIcon{2}{0}{1.4cm}};
  \node[mqttillustration] at (13.25,1.65) {\mqttIcon{0}{1}{1.4cm}};
  \node[mqttlabel,font=\scriptsize\bfseries] at (13.25,2.7) {Notificaciones};
  \node[mqttlabel,font=\scriptsize\bfseries] at (13.25,.85) {Auditoría};
  \draw[propuesta] (8.8,4.6) -- (8.8,4.0);
  \node[mqttdata,anchor=east] at (8.35,4.3) {SolicitudTutoriaCreada};
  \draw[propuesta] (9.9,2.95) -- (10.65,2.95) |- (12.4,3.5);
  \draw[propuesta] (10.65,2.95) |- (12.4,1.65);

  \node[mqttdata,text=turquesa,anchor=west] at (1.0,2.9) {sgra/tutorias/creada};
  \node[mqttlabel,anchor=west] at (1.0,2.2) {QoS 1\quad +\quad eventoId};
  \node[mqttdata,text=gris,anchor=west] at (1.0,1.5) {TLS / ACL\qquad Persistencia};
\end{tikzpicture}
\fuentes{\cite{sgra_frontend,sgra_backend,mqtt2019,shahri2022mqtt}}
\note{
La fila superior retoma la arquitectura de referencia documentada por el
equipo: interfaz Angular, backend Spring Boot y base de datos PostgreSQL.
La mensajería MQTT de la parte inferior es una propuesta de mejora;
las conexiones discontinuas indican los componentes y flujos propuestos.
Mosquitto es una posible implementación del broker, no una dependencia
que se haya demostrado como instalada en el SGRA.
\par
Después de registrar una solicitud, el backend podría publicar el evento
SolicitudTutoriaCreada. Notificaciones y auditoría podrían reaccionar
mediante suscripciones, sin una llamada directa del backend a cada consumidor.
Se podrían incorporar otros receptores mediante nuevas suscripciones.
\par
Como punto de partida se propone evaluar QoS 1 y un identificador único
del evento. Cada consumidor debe reconocer eventos ya procesados para
no repetir sus efectos cuando reciba duplicados. QoS 2 tampoco sustituye
la coordinación necesaria para ejecutar una operación de negocio una sola vez.
\par
Hay que resolver qué ocurre si el broker deja de estar disponible y qué
sucede si se confirma la transacción en la base de datos pero falla la
publicación del evento. MQTT no coordina por sí solo esas dos operaciones.
\par
TLS representa el cifrado del transporte; ACL, los permisos por tópico.
La persistencia y las sesiones deben configurarse según los requisitos
de recuperación. La demo inicial no proporciona una solución completa
de continuidad de servicio. QoS no establece un plazo máximo de entrega:
las afirmaciones de latencia y capacidad requieren mediciones en el sistema.
}
\end{frame}
\endgroup

% Sección reservada para nixon.
% Agregar aquí diapositivas propias con sus fuentes \cite{clave}.
% Mantener diagramas originales y evitar afirmar que la propuesta ya está implementada.
\begingroup
\section{JMS / Jakarta Messaging}

\begin{frame}{JMS / Jakarta Messaging: arquitectura}
{\large\bfseries API de mensajería para Java}\par\vspace{1mm}
{\small Un proveedor implementa la API; JMS no define un protocolo de red único.}
\begin{center}
\resizebox{0.91\textwidth}{!}{
\begin{tikzpicture}[x=1cm,y=1cm]
  \node[caja,text width=2.15cm,minimum height=1.2cm] (np) at (0,0)
    {\textbf{Productor Java}\\\scriptsize Backend SGRA};
  \node[evento,text width=2.65cm,minimum height=1.4cm] (nt) at (4.1,0)
    {\textbf{Proveedor compatible}\\\scriptsize Destino: Topic};
  \node[caja,text width=2.65cm] (nn) at (8.6,0.9)
    {\textbf{Notificaciones}\\\scriptsize Suscripción A};
  \node[caja,text width=2.65cm] (na) at (8.6,-0.9)
    {\textbf{Auditoría}\\\scriptsize Suscripción B};
  \draw[flecha] (np.east) -- node[above,font=\scriptsize]{publica} (nt.west);
  \draw[flecha] (nt.east) -- ++(0.65,0) |- (nn.west);
  \draw[flecha] (nt.east) -- ++(0.65,0) |- (na.west);
  \node[font=\scriptsize\ttfamily,text=gris] at (4.1,-1.65)
    {SolicitudTutoriaCreada};
\end{tikzpicture}
}
\end{center}
\begin{columns}[T,onlytextwidth]
\column{0.48\textwidth}
\textbf{Queue / cola}\par\vspace{1mm}
{\small Consumidores que compiten por los mensajes.}
\column{0.48\textwidth}
\textbf{Topic / tópico}\par\vspace{1mm}
{\small Cada suscripción independiente recibe el evento.}
\end{columns}
\fuentes{\cite{jakarta2022}}
\end{frame}

\begin{frame}{Ejemplo Java: un evento, dos suscripciones}
{\small\textbf{Práctica local:} Java + Apache Artemis 2.57.0.}
\vspace{2mm}
\begin{enumerate}
  \item \textbf{Terminal A:} notificaciones; \textbf{B:} auditoría.
  \item Ambas se suscriben a \texttt{sgra.tutorias.creada}.
  \item \textbf{Terminal C:} publica \texttt{SolicitudTutoriaCreada}.
\end{enumerate}
\begin{block}{Datos ficticios del mensaje}
\small
\texttt{eventoId: evt-001}\quad\texttt{solicitudId: 42}\\[1mm]
\texttt{tipo: SolicitudTutoriaCreada}
\end{block}
{\small\textbf{Resultado comprobado:} A y B recibieron el mismo evento\par
con el mismo \texttt{JMSMessageID}.}\par\vspace{2mm}
{\scriptsize\color{gris}Dos suscripciones no durables. Datos ficticios; sin conexión al SGRA.}
\fuentes{\cite{jakarta2022}}
\end{frame}

\begin{frame}{SGRA: propuesta de integración con JMS}
\begin{center}
\resizebox{0.68\textwidth}{!}{
\begin{tikzpicture}[x=1cm,y=1cm]
  \node[caja,minimum width=2.3cm] (ui) at (0,0) {\textbf{Angular}};
  \node[caja,minimum width=2.7cm] (api) at (4,0) {\textbf{Spring Boot}};
  \node[caja,minimum width=2.6cm] (db) at (8.3,0) {\textbf{PostgreSQL}};
  \draw[flecha] (ui) -- node[above,font=\scriptsize]{REST/HTTPS} (api);
  \draw[flecha] (api) -- node[above,font=\scriptsize]{SQL} (db);
  \node[evento,minimum width=2.7cm] (broker) at (4,-1.5)
    {\textbf{Proveedor JMS}};
  \node[caja,text width=2.6cm] (cons) at (8.3,-1.5)
    {\scriptsize Notificaciones\\Auditoría};
  \draw[flecha,dashed] (api) -- (broker);
  \draw[flecha,dashed] (broker) -- node[above,font=\scriptsize]{dos suscripciones} (cons);
\end{tikzpicture}
}
\end{center}
\begin{columns}[T,onlytextwidth]
\column{0.48\textwidth}
\textbf{Integración Java}
\begin{itemize}
  \item \texttt{JmsTemplate}: publicación.
  \item \texttt{@JmsListener}: recepción.
\end{itemize}
\textbf{Ventaja}\par\vspace{1mm}
{\small Añadir receptores sin llamadas directas desde el productor.}
\column{0.48\textwidth}
\textbf{Decisiones pendientes}
\begin{itemize}
  \item Guardar solicitud y publicar evento.
  \item Desconexiones y duplicados.
  \item Medir latencia y capacidad.
\end{itemize}
\end{columns}
\vspace{3mm}
\fuentes{\cite{springjms,maharjan2023}}
\end{frame}
\endgroup
% Sección de John Silva: AMQP/RabbitMQ y análisis comparativo, aplicada al SSDC.
% La propuesta es una evolución futura: el SSDC actual es un monolito y NO usa un broker.
\begingroup
\section{AMQP y análisis comparativo}

\begin{frame}{AMQP y RabbitMQ: exchange, binding y cola}
{\small El productor publica en un \textbf{exchange}; los \textbf{bindings} lo conectan con \textbf{colas} y los consumidores confirman (\textit{ack}).}
\begin{center}
\resizebox{0.9\textwidth}{!}{
\begin{tikzpicture}[x=1cm,y=1cm]
  \node[caja,text width=2.5cm,minimum height=1.3cm] (pub) at (0,0)
    {\textbf{Servicio de prepostulaciones}\\\scriptsize productor};
  \node[evento,text width=2.5cm,minimum height=1.3cm] (ex) at (4.3,0)
    {\textbf{Exchange topic}\\\scriptsize\texttt{ssdc.eventos}};
  \node[caja,text width=2.5cm] (q1) at (8.8,1.6)  {\scriptsize\texttt{q.correo}};
  \node[caja,text width=2.5cm] (q2) at (8.8,0)    {\scriptsize\texttt{q.notificaciones}};
  \node[caja,text width=2.5cm] (q3) at (8.8,-1.6) {\scriptsize\texttt{q.auditoria}};
  \node[caja,text width=2.3cm] (c1) at (12.6,1.6)  {\scriptsize Servicio de correo};
  \node[caja,text width=2.3cm] (c2) at (12.6,0)    {\scriptsize Servicio de notificaciones};
  \node[caja,text width=2.3cm] (c3) at (12.6,-1.6) {\scriptsize Servicio de auditoría};
  \draw[flecha] (pub) -- node[above,font=\scriptsize]{publica} (ex);
  \foreach \q in {q1,q2,q3}{\draw[flecha] (ex.east) -- ++(0.6,0) |- (\q.west);}
  \foreach \q/\c in {q1/c1,q2/c2,q3/c3}{\draw[flecha] (\q) -- (\c);}
  \node[font=\tiny\ttfamily,text=gris] at (6.65,2.35) {binding: prepostulacion.aprobada};
  \node[font=\scriptsize,text=gris] at (10.7,-2.45) {consumo con \textit{ack}};
\end{tikzpicture}}
\end{center}
{\small\textbf{Evento de ejemplo:} \texttt{PrepostulacionAprobada} (hoy, dentro del backend y sin broker).}\par\vspace{1mm}
{\scriptsize\color{gris}Matiz: exchange y binding son del modelo 0-9-1; AMQP 1.0 (OASIS) no los define.}
\fuentes{\cite{oasis2012,rabbitmq_concepts}}
\end{frame}

\begin{frame}{MQTT vs JMS vs AMQP: ¿cuál encaja en el SSDC?}
\small
\begin{center}
\renewcommand{\arraystretch}{1.25}
\resizebox{\textwidth}{!}{%
\begin{tabular}{@{}p{2.6cm}p{4.2cm}p{4.2cm}p{4.6cm}@{}}
\toprule
 & \textbf{MQTT} & \textbf{JMS} & \textbf{AMQP} \\
\midrule
Naturaleza & Protocolo ligero pub/sub & API Java; el proveedor define el protocolo & Protocolo de red estándar \\
Modelo & Tópicos & Colas y tópicos & Exchanges, bindings y colas \\
Entrega & QoS 0, 1 y 2 & Ack y mensajes persistentes & Ack, colas durables, confirmaciones \\
Encaje en el SSDC & Pensado para dispositivos con poco ancho de banda; el SSDC no los tiene & Integra bien con Spring, pero ata el diseño a Java & Independiente del lenguaje; permite enrutar a varios servicios \\
\bottomrule
\end{tabular}}
\end{center}
\begin{block}{Tolerancia a fallos del broker}
\small Si el broker cae, los servicios no se comunican: hace falta colas durables, \textit{ack} manual, cola de mensajes fallidos y un broker replicado.
\end{block}
\fuentes{\cite{derlatka2023,maharjan2023,bezerra2024,gultunca2018}}
\end{frame}

\begin{frame}{Recomendación para el SSDC y compromisos}
\footnotesize
\begin{columns}[T,onlytextwidth]
\column{0.52\textwidth}
\textbf{AMQP con RabbitMQ para eventos asíncronos}
\begin{itemize}
  \item Un exchange \texttt{topic} para todos los eventos.
  \item Colas durables y \textit{ack} manual.
  \item Cola de fallidos y reintentos acotados.
  \item Consumidores idempotentes (\texttt{eventoId}).
\end{itemize}
\column{0.44\textwidth}
\textbf{Compromisos}
\begin{itemize}
  \item Un componente más que operar y monitorear.
  \item Consistencia eventual: el correo llega después de aprobar.
  \item Entrega al menos una vez: pueden llegar duplicados.
\end{itemize}
\end{columns}
\begin{block}{Antes de adoptarlo}
\footnotesize Medir con la carga real del SSDC: los estudios comparativos no son mediciones de nuestro sistema \cite{maharjan2023}.
\end{block}
{\scriptsize\color{gris}Propuesta de evolución: el SSDC actual no usa broker.}
\fuentes{\cite{rabbitmq_concepts,maharjan2023}}
\end{frame}

\begin{frame}{Incidente real: Google Cloud, 2 de junio de 2019}
\small
\begin{center}
\resizebox{0.97\textwidth}{!}{
\begin{tikzpicture}[x=1cm,y=1cm]
  \node[caja,text width=3.2cm] (a) at (0,0)    {\textbf{Detonante}\\\scriptsize mantenimiento en una sola ubicación};
  \node[evento,text width=3.6cm] (b) at (4.6,0) {\textbf{Causa raíz}\\\scriptsize error de software en la automatización: detuvo clústeres de lugares distintos a la vez};
  \node[caja,text width=3.4cm] (c) at (9.4,0)  {\textbf{Defensa que falló}\\\scriptsize modo \textit{fail static} duró solo minutos; se retiraron rutas BGP};
  \node[caja,text width=3.0cm] (d) at (13.8,0) {\textbf{Efecto}\\\scriptsize pérdida de capacidad de red en seis regiones};
  \draw[flecha] (a) -- (b); \draw[flecha] (b) -- (c); \draw[flecha] (c) -- (d);
\end{tikzpicture}}
\end{center}
\begin{itemize}
  \item Duración: entre 3 h 19 min y 4 h 25 min según la región.
  \item Herramientas de depuración compitieron con la red congestionada.
\end{itemize}
\begin{block}{Qué aprendemos para el SSDC con broker}
\footnotesize Un cambio de configuración o de mantenimiento no debe aplicarse a todos los nodos del broker a la vez: se prueba primero en uno y se puede volver atrás.
\end{block}
\fuentes{\cite{googlecloud19009}}
\end{frame}
\endgroup


\renewcommand*{\bibfont}{\scriptsize} % tamaño legible para 16:9
% Para la entrega final: ampliar con las fuentes que citen los demás integrantes.
% Las seis revistas con DOI reciente se citan en las secciones y se imprimen.
\begin{frame}[allowframebreaks]{Referencias}
\printbibliography[heading=none]
\end{frame}
\end{document}
